% =============================================================================
%  Section 2.2 — Các hình thức biểu diễn tri thức
% =============================================================================

\section{Các hình thức biểu diễn tri thức}

Trong các hệ thống trí tuệ nhân tạo hiện đại, đặc biệt là các hệ thống khai thác tri thức có cấu trúc, Ontology và đồ thị tri thức giữ vai trò nền tảng trong việc tổ chức, chuẩn hóa và biểu diễn thông tin. Dữ liệu văn bản thuần túy chỉ phản ánh tri thức ở dạng phi cấu trúc, trong khi đó Ontology và đồ thị tri thức cho phép mô hình hóa rõ ràng các khái niệm, thực thể và quan hệ giữa chúng. Điều này có ý nghĩa lớn với các miền tri thức có cấu trúc quan hệ phức tạp như nghệ thuật Chèo, nơi tồn tại sự liên kết đa tầng giữa nhân vật, vở diễn, vai diễn, làn điệu, bối cảnh lịch sử, ... Phần này trình bày cơ sở lý thuyết của bản thể luận và đồ thị tri thức, đồng thời làm rõ vai trò của chúng trong việc xây dựng hệ thống truy xuất tri thức có cấu trúc.

\subsection{Bản thể luận - Ontology}

Ontology là một mô hình hình thức hóa tri thức của một miền xác định, trong đó các khái niệm và quan hệ giữa chúng được định nghĩa một cách tường minh và có cấu trúc. Trong khoa học máy tính, ontology là một hệ thống bao gồm các lớp khái niệm (classes), thuộc tính (properties) và các ràng buộc logic nhằm đảm bảo tính nhất quán ngữ nghĩa. Việc hình thức hóa này cho phép các mô hình ngôn ngữ lớn hiểu được cấu trúc của miền tri thức, thay vì chỉ xử lý dữ liệu ở mức ký hiệu hoặc chuỗi văn bản. Nhờ đó, ontology tạo tiền đề cho các cơ chế suy luận tự động và chuẩn hóa kho thông tin.

Cấu trúc của một ontology thường được tổ chức theo phân cấp khái niệm, trong đó các lớp có thể kế thừa thuộc tính và quan hệ từ lớp cha. Các thuộc tính được chia thành thuộc tính đối tượng (object properties), mô tả quan hệ giữa các thực thể, và thuộc tính dữ liệu (data properties), mô tả giá trị cụ thể của thực thể. Trong miền nghệ thuật Chèo, ontology có thể xác định rõ ràng mối quan hệ giữa các khái niệm như vở diễn, nhân vật, vai diễn và nghệ sĩ, từ đó tạo ra một khung khái niệm thống nhất cho toàn bộ hệ thống.

Về mặt kỹ thuật, ontology thường được biểu diễn bằng các chuẩn như Resource Description Framework (RDF), Resource Description Framework Schema (RDFS) và Web Ontology Language (OWL), trong đó OWL hỗ trợ biểu diễn các ràng buộc logic phức tạp và tích hợp với các bộ suy luận tự động. Trong các hệ thống biểu diễn tri thức hiện đại, ontology thường đóng vai trò như một tầng lược đồ khái niệm (schema layer), định nghĩa cấu trúc và ngữ nghĩa của miền tri thức. Dựa trên lược đồ này, các thực thể và quan hệ cụ thể có thể được tổ chức thành một đồ thị tri thức, cho phép lưu trữ và khai thác dữ liệu ở quy mô lớn.

\subsection{Đồ thị tri thức - Knowledge Graph}

Đồ thị tri thức là một mô hình biểu diễn tri thức dưới dạng đồ thị, trong đó các thực thể được biểu diễn bởi các nút (nodes) và các quan hệ giữa chúng được biểu diễn bởi các cạnh (edges). Mỗi đơn vị tri thức thường được mô tả dưới dạng bộ ba \textit{(subject – predicate – object)}, cho phép biểu diễn thông tin một cách linh hoạt, trực quan và dễ mở rộng. Khác với các mô hình cơ sở dữ liệu truyền thống vốn phụ thuộc vào cấu trúc bảng cố định, đồ thị tri thức cho phép bổ sung thực thể và quan hệ mới mà không làm thay đổi lược đồ tổng thể, từ đó phù hợp với các miền tri thức có tính mở và liên tục phát triển.

Trong hệ thống biểu diễn tri thức hiện đại, đồ thị tri thức được xây dựng dựa trên nền tảng Ontology. Ontology đóng vai trò là tầng khái niệm (schema layer), định nghĩa các lớp, thuộc tính và ràng buộc logic của miền tri thức, còn đồ thị tri thức là tầng dữ liệu (instance layer), trong đó các thực thể cụ thể và các mối quan hệ giữa chúng được hiện thực hóa theo cấu trúc đã được quy định. Sự phân tách này giúp đảm bảo rằng dữ liệu trong đồ thị tri thức được tổ chức một cách có hệ thống, đồng thời tuân thủ các ràng buộc ngữ nghĩa, từ đó hỗ trợ khả năng suy luận và kiểm tra tính nhất quán.

Về mặt triển khai, đồ thị tri thức có thể được xây dựng theo hai hướng tiếp cận chính. Thứ nhất là mô hình RDF triple store, trong đó dữ liệu được lưu trữ dưới dạng các bộ ba và phù hợp với các hệ thống chú trọng suy luận ngữ nghĩa dựa trên các chuẩn như RDF, RDFS và OWL. Thứ hai là mô hình property graph, được sử dụng trong các hệ quản trị cơ sở dữ liệu đồ thị như Neo4j, cho phép gán thuộc tính trực tiếp lên các nút và cạnh, từ đó hỗ trợ hiệu quả các truy vấn theo đường đi và truy vấn đa bước. Việc lựa chọn mô hình phụ thuộc vào yêu cầu cụ thể của hệ thống, bao gồm khả năng suy luận, hiệu năng truy vấn và tính linh hoạt trong mở rộng dữ liệu.

Đồ thị tri thức có khả năng biểu diễn các quan hệ đa chiều và phức tạp mà không làm gia tăng độ phức tạp của lược đồ. Đồ thị tri thức không cần thực hiện các phép nối (join) phức tạp như trong cơ sở dữ liệu quan hệ, các truy vấn trên đồ thị tri thức có thể khai thác trực tiếp cấu trúc liên kết giữa các thực thể thông qua các đường đi trong đồ thị. Điều này đặc biệt hữu ích trong các bài toán yêu cầu khám phá tri thức hoặc truy vấn theo ngữ cảnh, nơi mà các mối quan hệ giữa dữ liệu đóng vai trò quan trọng. Ngoài ra, đồ thị tri thức còn đóng vai trò quan trọng trong việc nâng cao khả năng giải thích của hệ thống. Mỗi kết quả truy vấn không chỉ là một giá trị đầu ra, mà còn có thể được truy vết thông qua chuỗi các quan hệ cụ thể giữa các thực thể trong đồ thị. Điều này giúp tăng tính minh bạch và độ tin cậy, đặc biệt trong các hệ thống hỏi đáp yêu cầu khả năng giải thích quyết định.

Trong bối cảnh các hệ thống truy hồi tri thức hiện đại, đặc biệt là Retrieval-Augmented Generation (RAG), đồ thị tri thức được sử dụng như một nguồn tri thức có cấu trúc nhằm hỗ trợ quá trình truy vấn và mở rộng ngữ cảnh \cite{fan2023cichmkg,wan2024wumkg}. Khi được kết hợp với Ontology, đồ thị tri thức không chỉ cung cấp dữ liệu mà còn kế thừa các ràng buộc ngữ nghĩa của miền tri thức, từ đó nâng cao chất lượng truy hồi và hỗ trợ các cơ chế suy luận. Trong phạm vi nghiên cứu này, đồ thị tri thức đóng vai trò là lớp lưu trữ và khai thác dữ liệu, được xây dựng dựa trên Ontology nhằm đảm bảo tính nhất quán và phục vụ các thành phần khác của hệ thống.
